% Équivalent LaTeX de S02-verso-tableau.html (feuille élève, verso — tableau
% de consignation des réponses J1 à J4). Document autonome, distinct du
% livret LaTeX complet (voir README.md du dossier "Automatismes Hebdo") :
% se compile seul, sans dépendre de CLT2627.sty ni du pipeline
% AutomatismesSeconde2627.tex. Compiler avec pdflatex ou lualatex.
\documentclass[10pt]{article}
% Page 1 (recto) en portrait, page 2 (verso, le tableau) reste en paysage :
% geometry pose le format par défaut (portrait) et pdflscape fournit
% l'environnement landscape qui fait pivoter uniquement le verso (MediaBox
% de cette page seule tournée à 90°, indépendamment du reste du document).
\usepackage[a4paper,top=5mm,bottom=5mm,left=10mm,right=10mm]{geometry}
\usepackage{pdflscape}
\usepackage[french]{babel}
\usepackage[table]{xcolor}
\usepackage{array}
\usepackage{multirow}
\usepackage{pifont}
\usepackage[most]{tcolorbox}
\tcbuselibrary{raster}
\usepackage{ragged2e}
\usepackage{amssymb}
\usepackage{fontawesome5}
\usepackage{enumitem}
\setlist[itemize]{topsep=1pt, itemsep=1pt, parsep=0pt, leftmargin=1.3em, label=$\star$}
\usepackage{bbding}
\usepackage{multicol}
\usepackage{tikz}
\usetikzlibrary{calc}

\pagestyle{empty}
\setlength{\parindent}{0pt}
\renewcommand{\familydefault}{\sfdefault}

% --- Palette (reprise des tokens CSS de S02-verso-tableau.html) -----------
\definecolor{inkC}{HTML}{20232B}
\definecolor{inksoftC}{HTML}{565A66}
\definecolor{mutedC}{HTML}{8B8F99}
\definecolor{rulestrongC}{HTML}{B9C2D1}
\definecolor{qcmC}{HTML}{2455B8}
\definecolor{questionC}{HTML}{12805C}
\definecolor{finalC}{HTML}{B5511F}
\definecolor{bonusC}{HTML}{B8298F}
\definecolor{bonussoftC}{HTML}{FBE6F3}
\definecolor{emptyC}{HTML}{E7E9EE}
\definecolor{redC}{HTML}{B8342B}
% Couleurs des boîtes "À retenir"/"Exemple", reprises telles quelles de
% \ARetenirBox/\ExempleBoxAuto dans CLT2627.sty (bleuPropriete, grisExemple)
% pour rester visuellement fidèle au fichier original.
\definecolor{bleuPropriete}{RGB}{41,121,204}
\definecolor{grisExemple}{RGB}{85,85,95}

\newcommand{\starQ}{\textcolor{questionC}{$\bigstar$}}
\newcommand{\starF}{\textcolor{finalC}{$\bigstar$}}
\newcommand{\starFF}{\textcolor{finalC}{$\bigstar\bigstar$}}
\newcommand{\okmark}{\ding{51}}
\newcommand{\nomark}{\ding{55}}
% Case à cocher vierge (l'élève y trace lui-même un crochet ou une croix).
\newcommand{\corrbox}{\fboxsep=0pt\fbox{\rule{0pt}{8pt}\rule{8pt}{0pt}}}

% Case active : numéro de question entouré en haut à gauche de la case,
% le reste de la case restant vide pour l'écriture manuscrite. (Après
% correction collective, l'élève entoure ou coche ce numéro pour indiquer
% juste/faux -- voir légende.)
\newcommand{\qcase}[1]{%
  \begin{tikzpicture}[baseline=0pt]
    \useasboundingbox (0,0) rectangle (2.95cm,1.2cm);
    \node[anchor=north west, circle, draw=mutedC, line width=0.4pt,
          inner sep=1pt, minimum size=1.1em, font=\tiny, text=mutedC]
      at (0.08cm,1.12cm) {#1};
  \end{tikzpicture}%
}
% Case grisée : aucune question ce jour-là pour cet automatisme.
\newcommand{\qempty}{\cellcolor{emptyC}}

% Bandeau de titre du recto : petite étiquette en haut à gauche qui entaille
% le cadre, grand titre centré en dessous. Repris du style \FeuilletTitre
% utilisé dans les feuilles d'exercices des séquences (même famille visuelle),
% reconstruit ici pour rester autonome (pas de dépendance à CLT2627.sty).
\NewDocumentCommand{\RectoTitre}{m m}{%
  \begin{center}
  \begin{tcolorbox}[
    enhanced,
    width=0.8\linewidth,
    colback=white,
    colframe=inkC,
    boxrule=0.8pt,
    arc=8pt,
    sharp corners=downhill,
    left=10pt, right=10pt, top=4pt, bottom=3pt,
    halign=center,
    overlay={
      \node[fill=white, anchor=west, font=\normalsize, inner xsep=4pt]
        at ([xshift=14pt]frame.north west) {#1};
    },
  ]
  {\huge\bfseries #2}
  \end{tcolorbox}
  \end{center}
}

% Carte "automatisme" du recto : petit engrenage (repris de \AutoItem dans
% AutomatismesSeconde2627.tex, où \faCogs symbolise l'automatisme) + code
% + intitulé, tout sur une seule ligne (\tcbox se dimensionne à son contenu,
% pas de retour à la ligne forcé), même code couleur que l'en-tête du
% tableau du verso (qcmC / questionC / finalC).
\newcommand{\autoCardLine}[3]{%
  \noindent\tcbox[
    colback=white, colframe=#1, coltext=#1,
    boxrule=1.2pt, arc=5pt, left=10pt, right=10pt, top=2pt, bottom=2pt,
  ]{{\color{#1}\faCogs}~\textcolor{#1}{\Large\bfseries #2}~~{\normalsize\color{inksoftC} #3}}%
  \par
}

% Boîte "À retenir" : même cadre que les boîtes "Propriété" du cours
% (\Prp dans CLT2627.sty) -- pas de bordure pleine, seulement 4 coins en
% équerre dessinés en overlay, fond bleuPropriete!6.
\newcommand{\ARetenirBox}[1]{%
  \begin{tcolorbox}[
    enhanced,
    colback=bleuPropriete!6, colframe=bleuPropriete,
    boxrule=0pt, arc=0pt, outer arc=0pt,
    left=10pt, right=8pt, top=2pt, bottom=1pt,
    overlay={
      \draw[bleuPropriete, line width=1.2pt] (frame.north west) -- ++(7pt,0pt);
      \draw[bleuPropriete, line width=1.2pt] (frame.north west) -- ++(0pt,-7pt);
      \draw[bleuPropriete, line width=1.2pt] (frame.north east) -- ++(-7pt,0pt);
      \draw[bleuPropriete, line width=1.2pt] (frame.north east) -- ++(0pt,-7pt);
      \draw[bleuPropriete, line width=1.2pt] (frame.south west) -- ++(7pt,0pt);
      \draw[bleuPropriete, line width=1.2pt] (frame.south west) -- ++(0pt,7pt);
      \draw[bleuPropriete, line width=1.2pt] (frame.south east) -- ++(-7pt,0pt);
      \draw[bleuPropriete, line width=1.2pt] (frame.south east) -- ++(0pt,7pt);
    },
  ]
    \textbf{\textcolor{bleuPropriete}{À retenir}}\par\vspace{1pt}
    #1
  \end{tcolorbox}%
}
% Boîte "Exemple" : même liseré que \ExempleBoxAuto, plus le crayon
% encerclé en overlay au coin supérieur gauche, repris à l'identique de
% \Exemple dans CLT2627.sty (\PencilRightUp, police bbding).
\newcommand{\ExempleBoxR}[1]{%
  \begin{tcolorbox}[
    enhanced,
    colback=grisExemple!6, colframe=grisExemple,
    boxrule=0pt, borderline west={2pt}{0pt}{grisExemple},
    arc=0pt, outer arc=0pt,
    left=16pt, right=8pt, top=3pt, bottom=1pt,
    overlay={
      \node[anchor=center, circle, fill=white, draw=grisExemple, line width=0.8pt, minimum size=18pt] at ([yshift=-10pt]frame.north west) {\textcolor{grisExemple}{\fontsize{11}{11}\selectfont\PencilRightUp}};
    },
  ]
    \textbf{\textcolor{grisExemple}{Exemple}}\par\vspace{1pt}
    #1
  \end{tcolorbox}%
}
% Place une \ARetenirBox et une \ExempleBoxR côte à côte, chacune à sa
% propre hauteur : pas de "raster equal height", pour que la boîte
% "À retenir" ne s'étire pas inutilement jusqu'à la hauteur, souvent plus
% grande, de la boîte "Exemple" en vis-à-vis.
\newcommand{\RecapPair}[2]{%
  \begin{tcbraster}[raster columns=2, raster column skip=8pt, raster force size=false, raster valign=top, before skip=0pt, after skip=0pt]
    \ARetenirBox{\normalsize#1}
    \ExempleBoxR{\footnotesize#2}
  \end{tcbraster}%
}

\newcolumntype{J}{>{\RaggedRight\arraybackslash}m{1.7cm}}
\newcolumntype{N}{>{\centering\arraybackslash}m{1.05cm}}
\newcolumntype{O}{>{\centering\arraybackslash}m{3.15cm}}

\begin{document}

% ============================= RECTO (page 1) =============================
\RectoTitre{ Semaine du 07 au 11 septembre}{Automatismes Série 1}

%\vspace{0pt}
%\begin{center}
%{\large\bfseries\color{inkC}Cette semaine, 7 automatismes sont abordés en questions rituelles.}
%\end{center}

\vspace{0pt}

\autoCardLine{qcmC}{A1}{Comparer deux nombres}
\vspace{1pt}
\RecapPair{%
  Pour comparer des fractions, on peut~:
  \begin{itemize}
    \item les réduire au même dénominateur \textbf{positif} et comparer
    leurs numérateurs~: les fractions sont dans le même ordre que leurs
    numérateurs respectifs~;
    \item les comparer à~$1$.
  \end{itemize}
}{%
  Pour comparer les fractions $\dfrac56$ et $\dfrac34$, on cherche un
  dénominateur commun à ces deux fractions. Le plus petit multiple commun
  aux nombres $6$ et $4$ est $12$.
  \[\dfrac56=\dfrac{5\times2}{6\times2}=\dfrac{10}{12} \quad\text{et}\quad
  \dfrac34=\dfrac{3\times3}{4\times3}=\dfrac{9}{12}\]
  On en déduit que $\dfrac{10}{12}>\dfrac{9}{12}$. Ainsi, $\dfrac56>\dfrac34$.
}
\vspace{1pt}
\RecapPair{%
  $\star$ Pour comparer deux nombres réels $A$ et $B$, on peut~:
  \begin{itemize}
    \item étudier le signe de leur différence. Par exemple, si $A-B>0$,
    alors $A>B$~;
    \item comparer leur quotient à~$1$, s'ils sont strictement positifs. Par
    exemple, si $\dfrac{A}{B}<1$, alors $A<B$.
  \end{itemize}
}{%
  $\star$ Pour comparer $A=\dfrac45$ et $B=\dfrac56$, on calcule la différence
  $A-B$~:
  \[A-B=\dfrac45-\dfrac56=\dfrac{4\times6}{5\times6}-\dfrac{5\times5}{6\times5}=\dfrac{24}{30}-\dfrac{25}{30}=-\dfrac1{30}\]
  Comme $-\dfrac1{30}<0$, on a $A-B<0$. Par conséquent $A<B$.

  \vspace{0pt}
  Pour comparer $A=3^{60}\times2^{62}$ et $B=6^{61}$, on peut comparer le
  quotient $\dfrac AB$ avec $1$ car les deux nombres sont strictement
  positifs.
  \[\dfrac AB=\dfrac{3^{60}\times2^{62}}{6^{61}}=\dfrac{3^{60}\times2^{62}}{(3\times2)^{61}}=\dfrac{3^{60}\times2^{62}}{3^{61}\times2^{61}}=2^{1}\times3^{-1}=\dfrac23\]
  Comme $\dfrac AB<1$, on a~: $A<B$.
}
\vspace{0pt}

\vspace*{0.12cm}

\begin{multicols}{3}
\autoCardLine{qcmC}{A2}{Fractions}
\columnbreak
\autoCardLine{qcmC}{A3}{Puissances}
\columnbreak
\textit{Voir Cartes Mentales et Cours Séquence~1.}
\end{multicols}
\vspace{0pt}

\autoCardLine{qcmC}{A4}{Écritures d'un nombre}
\vspace{1pt}
\RecapPair{%
  Tout nombre décimal peut s'écrire comme la somme de sa partie entière et
  de fractions décimales (dixièmes~: $\frac1{10}$, centièmes~:
  $\frac1{100}$, millièmes~: $\frac1{1\,000}$, etc.). \\
  Il peut aussi
  s'écrire sous la forme d'une seule fraction décimale.
}{%
  $4,5=4+\dfrac5{10}=\dfrac{45}{10}$

  $34,52=34+\dfrac5{10}+\dfrac2{100}=\dfrac{3\,452}{100}$

  $0,125=\dfrac1{10}+\dfrac2{100}+\dfrac5{1\,000}=\dfrac{125}{1\,000}$
}
\vspace{1pt}
\RecapPair{%
  La notation scientifique est l'écriture sous la forme $a\times10^n$, où
  $a$ possède un seul chiffre non nul avant la virgule \\
  ($1\leqslant a<10$)
  et $n$ est un entier relatif.
}{%
  Grand nombre~: $5\,230=5,23\times1\,000=5,23\times10^3$

  Petit nombre~: $0,007\,4=7,4\times0,001=7,4\times10^{-3}$
}

\vspace*{0.12cm}

\autoCardLine{questionC}{A22}{Pythagore}
\vspace{1pt}
\RecapPair{%
  Si $ABC$ est rectangle en $A$, alors $BC^2=AB^2+AC^2$ ($BC$ hypoténuse,
  côté opposé à l'angle droit).
}{%
  \begin{minipage}[c]{0.36\linewidth}
  \centering
  \begin{tikzpicture}[scale=1, line width=0.6pt]
    \coordinate (A) at (0,0);
    \coordinate (B) at (0,1);
    \coordinate (C) at (2.2,0);
    \draw (A) -- (B) -- (C) -- cycle;
    \draw (0,0.2) -- (0.2,0.2) -- (0.2,0);
    \node[left=2pt] at ($(A)!0.5!(B)$) {$4$};
    \node[below=2pt] at ($(A)!0.5!(C)$) {$8$};
    \node[above right=-2pt] at ($(B)!0.5!(C)$) {$?$};
    \node[below left=1pt] at (A) {$A$};
    \node[above left=1pt] at (B) {$B$};
    \node[below right=1pt] at (C) {$C$};
  \end{tikzpicture}
  \end{minipage}%
  \hfill
  \begin{minipage}[c]{0.6\linewidth}
  $ABC$ est rectangle en $A$, donc d'après le th\'eor\`eme de Pythagore~: \\
  $BC^2=AB^2+AC^2$.\\[2pt]
  Avec $AB=4$ et $AC=8$~: $BC^2=4^2+8^2=16+64=80$, \\
  donc $BC=\sqrt{80}=4\sqrt5$.
  \end{minipage}
}
\vspace{0pt}

\begin{multicols}{3}
\autoCardLine{finalC}{Algo}{Lire un programme}
\columnbreak
\autoCardLine{finalC}{Translations}{}
\columnbreak
\textit{Pas de rappels. \\
{\small Les questions ne nécessitent pas de pré-requis.}}
\end{multicols}

\newpage
% ============================= VERSO (page 2) ==============================
\begin{landscape}
\noindent
Nom~: \hrulefill{} \hspace{14pt} Prénom~: \hrulefill{} \hspace{14pt} Classe~: \hrulefill

\vspace{4pt}
\footnotesize
\textcolor{qcmC}{\textbf{QCM}} \hspace{16pt}
\starQ{} = Question \hspace{16pt}
\starFF{} = Automatisme final \hspace{16pt}
\colorbox{emptyC}{\rule{0pt}{5pt}\hspace{10pt}} pas de question ce jour \hspace{16pt}
\textcolor{questionC}{\okmark} juste \textperiodcentered{} \textcolor{redC}{\nomark} faux

\vspace{3pt}
\renewcommand{\arraystretch}{0.95}
\setlength{\tabcolsep}{2.5pt}
\begin{center}
\begin{tabular}{|J|N||*{7}{O|}}
\hline
\multicolumn{2}{|c||}{} &
\shortstack{\textcolor{qcmC}{A1}\\{\tiny Comparer deux nombres}} &
\shortstack{\textcolor{qcmC}{A2}\\{\tiny Fractions}} &
\shortstack{\textcolor{qcmC}{A3}\\{\tiny Puissances}} &
\shortstack{\textcolor{qcmC}{A4}\\{\tiny Écritures d'un nombre}} &
\shortstack{\textcolor{questionC}{A22}\\{\tiny Pythagore, Thalès}} &
\shortstack{\textcolor{finalC}{Algo}\\{\tiny Lire un programme}} &
\shortstack{\textcolor{finalC}{Translations}\\{\tiny (non numéroté)}} \\
\hline\hline

\multirow{3}{1.6cm}{\textbf{Jour 1}\\[-1pt] Date~:\\[-1pt] \hrulefill}
  & \textcolor{qcmC}{\textbf{QCM}}
  & \qcase{1} & \qcase{2} & \qempty & \qempty & \qempty & \qempty & \qempty \\
\cline{2-9}
  & \starQ
  & \qempty & \qcase{3} & \qcase{4} & \qempty & \qempty & \qempty & \qempty \\
\cline{2-9}
  & \starFF
  & \qempty & \qcase{5} & \qcase{6} & \qempty & \qempty & \qempty & \qempty \\
\hline

\multirow{3}{1.6cm}{\textbf{Jour 2}\\[-1pt] Date~:\\[-1pt] \hrulefill}
  & \textcolor{qcmC}{\textbf{QCM}}
  & \qempty & \qempty & \qempty & \qcase{1} & \qcase{2} & \qempty & \qempty \\
\cline{2-9}
  & \starQ
  & \qcase{3} & \qempty & \qempty & \qcase{4} & \qempty & \qempty & \qempty \\
\cline{2-9}
  & \starFF
  & \qcase{5} & \qempty & \qempty & \qempty & \qcase{6} & \qempty & \qempty \\
\hline

\multirow{3}{1.6cm}{\textbf{Jour 3}\\[-1pt] Date~:\\[-1pt] \hrulefill}
  & \textcolor{qcmC}{\textbf{QCM}}
  & \qempty & \qempty & \qempty & \qempty & \qempty & \qcase{1} & \qcase{2} \\
\cline{2-9}
  & \starQ
  & \qempty & \qempty & \qempty & \qempty & \qempty & \qcase{3} & \qcase{4} \\
\cline{2-9}
  & \starFF
  & \qcase{5} & \qcase{6} & \qempty & \qempty & \qempty & \qempty & \qempty \\
\hline

\multirow{3}{1.6cm}{\textbf{Jour 4}\\[-1pt] Date~:\\[-1pt] \hrulefill}
  & \textcolor{qcmC}{\textbf{QCM}}
  & \qempty & \qempty & \qempty & \qcase{1} & \qempty & \qcase{2} & \qempty \\
\cline{2-9}
  & \starQ
  & \qcase{3} & \qempty & \qcase{4} & \qempty & \qempty & \qempty & \qempty \\
\cline{2-9}
  & \starFF
  & \qempty & \qempty & \qcase{5} & \qempty & \qcase{6} & \qempty & \qempty \\
\hline
\end{tabular}
\end{center}

\vspace{4pt}
\begin{tcolorbox}[
  boxrule=1.2pt, colframe=bonusC, colback=bonussoftC, arc=3pt,
  left=10pt, right=10pt, top=3pt, bottom=3pt
]
\textcolor{bonusC}{\textbf{\footnotesize Q7 $\cdot$ BONUS $\cdot$ JOUR 4}} \hspace{10pt}
{\footnotesize\color{inksoftC} A1 --- Comparer $a=6^{53}\times7^{55}$ et $b=42^{54}$ (question de
dépassement, hors des 6 questions du jour)} \hfill
\rule{2.6cm}{0.4pt}\hspace{6pt}\corrbox
\end{tcolorbox}
\end{landscape}

\end{document}
